\section{Технический проект}
\subsection{Обоснование выбора технологий проектирования}
\subsubsection{Язык программирования C\#}
В качестве основного языка программирования для реализации игрового движка и редактора сцен выбран язык C\#. Данный выбор обусловлен совокупностью факторов, включающих производительность, безопасность памяти, скорость разработки и доступность инструментальной поддержки. C\# является объектно-ориентированным языком с управляемой средой выполнения, что обеспечивает автоматическое управление памятью через сборщик мусора. Это позволяет снизить риск ошибок, связанных с утечками памяти и висячими указателями, которые характерны для языков с ручным управлением памятью, таких как C++. В контексте разработки игрового движка, где надёжность и долговременная стабильность работы имеют большое значение, автоматическое управление памятью существенно упрощает разработку и отладку.

Ключевым преимуществом C\# является богатая стандартная библиотека и экосистема .NET. Стандартная библиотека предоставляет удобные коллекции, средства работы с файловой системой, потоками и асинхронным программированием. Механизмы рефлексии позволяют реализовать гибкую систему сериализации сцен и компонентов без необходимости написания шаблонного кода. Система LINQ упрощает выполнение операций над коллекциями игровых объектов при реализации игровой логики. Кроме того, среда выполнения .NET включает средства низкоуровневого взаимодействия с нативным кодом через механизм P/Invoke, что необходимо для вызова функций OpenGL и управления окнами через библиотеку GLFW.

Важным аспектом выбора C\# является его производительность, достаточная для разработки игровых движков среднего уровня сложности. Благодаря Just-In-Time компиляции в сочетании с оптимизациями на уровне среды выполнения, код на C\# может достигать производительности, приближающейся к нативному коду на C++. При этом современные версии .NET (6.0 и выше) включают аппаратные векторы, интринсики и возможности low-level оптимизаций, позволяющие эффективно работать с массивами вершинных данных и матрицами преобразований. Критически важные для производительности участки, такие как обработка вершинных буферов и расчёт матриц трансформации, могут быть реализованы с использованием типов-структур, что позволяет избежать аллокаций в управляемой куче и снизить нагрузку на сборщик мусора.

Для разработки игрового движка также важна кроссплатформенность языка C\#. Хотя в рамках данной работы основной целевой платформой является Windows, код движка спроектирован таким образом, что его можно скомпилировать и запустить под управлением Linux или macOS с минимальными изменениями. Это обеспечивается использованием кроссплатформенных библиотек (GLFW, OpenGL) и независимостью от платформозависимых API Windows. Среда выполнения .NET доступна на всех основных операционных системах благодаря проекту .NET Core, что даёт возможность портирования движка на другие платформы при необходимости.

Наконец, выбор C\# обусловлен также удобством разработки графического интерфейса редактора сцен с использованием Windows Forms. Данная технология является зрелой, хорошо документированной и тесно интегрированной с C\#, что позволяет быстро создавать функциональные окна, панели инструментов, текстовые поля и списки с минимальными затратами времени. Использование единого языка программирования как для движка, так и для редактора сцен упрощает поддержку и развитие проекта, поскольку разработчику не требуется переключаться между разными языками и инструментами. Таким образом, язык C\# обеспечивает оптимальный баланс между производительностью, надёжностью и скоростью разработки для целей дипломного проекта.

\subsubsection{Графический API OpenGL}
В качестве основного графического интерфейса приложения для реализации рендеринга выбран OpenGL версии 4.6. Данный выбор обусловлен совокупностью факторов, включающих кроссплатформенность, достаточный уровень контроля над графическим конвейером, относительную простоту освоения и наличие обширной документации. OpenGL является открытым стандартом, разрабатываемым консорциумом Khronos Group, и поддерживается на всех современных операционных системах, включая Windows, Linux и macOS, что позволяет в перспективе портировать разрабатываемый движок на другие платформы без замены графического ядра.

Основным преимуществом OpenGL по сравнению с альтернативными API, такими как Vulkan или DirectX 12, является умеренная сложность реализации. Для создания минимального работающего приложения на OpenGL требуется около двухсот-трёхсот строк кода, в то время как для Vulkan аналогичный объём функциональности потребовал бы нескольких тысяч строк только для инициализации. В контексте дипломной работы с ограниченными временными ресурсами это позволяет сосредоточиться на реализации непосредственной функциональности игрового движка, а не на решении низкоуровневых проблем управления синхронизацией и распределения памяти. OpenGL при этом сохраняет достаточный контроль над графическим конвейером для реализации трёхмерной сцены с источниками света, текстурами и трансформациями объектов.

Версия OpenGL 4.6 была выбрана как минимальная целевая версия по следующим причинам. Данная версия включает Direct State Access, позволяющий изменять состояние объектов без привязки к контексту, что упрощает многопоточную работу с ресурсами. Поддерживаются шейдерные хранимые буферы, предоставляющие возможность эффективно передавать большие объёмы данных в шейдеры. Вычислительные шейдеры позволяют выполнять общие вычисления на графическом процессоре без привязки к графическому конвейеру. Кроме того, OpenGL 4.6 является современным стандартом, поддерживаемым большинством видеокарт, выпущенных после 2017 года, включая все дискретные видеокарты NVIDIA серии GeForce 900 и новее, а также AMD Radeon серии R9 300 и новее. Это обеспечивает достаточную совместимость с оборудованием, доступным в учебных аудиториях и у потенциальных пользователей.

Для взаимодействия с OpenGL из управляемого кода C\# используется механизм P/Invoke, позволяющий вызывать нативные функции, экспортируемые динамическими библиотеками драйвера видеокарты. Для упрощения работы с функциями OpenGL используется библиотека OpenTK, предоставляющая корректные сигнатуры всех функций OpenGL, а также типы данных, совместимые с GLSL по именованию и структуре. Альтернативой могло бы быть использование библиотеки Silk.NET, однако OpenTK выбрана благодаря её большей распространённости и стабильности. Критически важные вызовы OpenGL, такие как отрисовка буферов и смена шейдерных программ, выполняются часто, и каждый такой вызов является переходом через границу управляемого и неуправляемого кода. Для снижения накладных расходов на межъязыковые вызовы в движке используется пакетирование операций: множество вершинных данных передаётся в GPU одним вызовом, а не серией мелких операций.

Архитектурно работа с OpenGL в движке организована по принципу инкапсуляции. Классы, представляющие графические ресурсы, содержат идентификаторы объектов OpenGL и реализуют интерфейс IDisposable для гарантированного освобождения неуправляемых ресурсов. При удалении игрового объекта или выгрузке сцены соответствующие вызовы glDeleteBuffers, glDeleteTextures и glDeleteProgram выполняются автоматически. Весь код, связанный с OpenGL, изолирован в отдельном модуле графической подсистемы, что позволяет в будущем заменить OpenGL на другой графический API (например, Vulkan) без изменения остальных компонентов движка. Таким образом, выбор OpenGL 4.6 обеспечивает баланс между функциональностью, совместимостью, производительностью и сложностью разработки, полностью соответствующий целям и ограничениям дипломного проекта.

\subsubsection{Windows Forms для реализации редактора сцен}
Для реализации графического интерфейса редактора сцен выбрана технология Windows Forms. Данный выбор обусловлен простотой разработки, нативной интеграцией с операционной системой Windows и тесной совместимостью с языком C\#, который используется для разработки основного движка. Windows Forms является зрелой и хорошо документированной технологией, предоставляющей разработчику богатую коллекцию стандартных элементов управления, таких как кнопки, текстовые поля, списки, чекбоксы и панели, что позволяет быстро создавать функциональный интерфейс без необходимости написания сложного кода для отрисовки элементов.

Основным преимуществом Windows Forms для целей данной работы является простота интеграции с трёхмерной визуализацией на основе OpenGL. Windows Forms предоставляет элемент Panel, внутрь которого можно встроить окно OpenGL через получение дескриптора окна (HWND) и передачу его в библиотеку GLFW или непосредственно в OpenGL. Это позволяет размещать область трёхмерного предпросмотра сцены непосредственно внутри окна редактора, наряду с панелями инструментов и элементами управления. Альтернативные технологии, такие как WPF, предоставляют более современные возможности кастомизации интерфейса, но требуют значительно больше усилий для интеграции с нативным рендерингом OpenGL из-за особенностей своей композиции и использования DirectX. В контексте дипломной работы, где приоритетом является создание работающего прототипа, а не излишне сложного интерфейса, Windows Forms является наиболее прагматичным выбором.

Архитектурно редактор сцен на Windows Forms представляет собой главное окно, содержащее следующие основные элементы управления. Панель трёхмерного предпросмотра отображает текущую сцену в режиме реального времени, используя ту же графическую подсистему OpenGL, что и игровой движок. Кнопка импорта модели открывает стандартный диалог выбора файла и запускает процесс загрузки STL или GLB модели. Текстовые поля для настройки камеры редактора позволяют изменять чувствительность управления, скорость перемещения, ближнюю и дальнюю дистанции отсечения. Список превью отображает все загруженные в сцену игровые объекты в виде миниатюр или текстовых названий, обеспечивая быстрый доступ к любому объекту для его выбора и редактирования. Отдельное окно инспектора объекта открывается при выборе объекта в сцене или в списке превью и содержит поля для редактирования трансформации объекта, чекбокс видимости в игровом режиме, поле для задания тега, а также кнопки удаления объекта и возврата к главному окну.

Важной особенностью реализации является разделение режимов работы редактора. В режиме редактирования пользователь может свободно перемещать камеру, выбирать объекты и изменять их параметры. В режиме игры редактор скрывает элементы управления и передаёт управление игровой камере, параметры которой были настроены разработчиком. Переключение между режимами осуществляется по нажатию соответствующей кнопки или сочетания клавиш. Такое разделение позволяет разработчику игр тестировать сцену в условиях, приближенных к реальному игровому процессу, не выходя из редактора.

Обработка событий в Windows Forms организована через стандартную событийную модель. При нажатии кнопки импорта модели вызывается соответствующий обработчик, который запускает процедуру загрузки файла и добавления объекта в сцену. При изменении значения в текстовом поле обновляется соответствующий параметр камеры или трансформации объекта. При клике по объекту в трёхмерной области предпросмотра выполняется ray casting для определения выбранного объекта и открывается инспектор. Все события обрабатываются в основном потоке пользовательского интерфейса, а ресурсоёмкие операции, такие как импорт больших моделей, выносятся в асинхронные задачи с отображением индикатора загрузки, чтобы не блокировать интерфейс.

Таким образом, выбор Windows Forms для реализации редактора сцен обусловлен оптимальным соотношением скорости разработки, функциональности и простоты интеграции с OpenGL-рендерингом. Технология позволяет создать полноценный инструмент для работы со сценой, не отвлекаясь на низкоуровневую реализацию элементов управления, что особенно важно в рамках дипломного проекта с ограниченными временными ресурсами.

\subsubsection{Библиотека SixLabors.ImageSharp}
Для загрузки и обработки текстурных изображений в разрабатываемом игровом движке используется библиотека SixLabors.ImageSharp. Данная библиотека представляет собой полностью управляемую, кроссплатформенную и высокопроизводительную библиотеку для работы с двухмерной графикой на языке C\#. ImageSharp поддерживает широкий спектр форматов изображений, включая PNG, JPEG, BMP, GIF, TIFF, а также современный формат WebP. В отличие от нативных решений, библиотека не имеет внешних зависимостей и полностью написана на управляемом коде, что упрощает распространение приложения и исключает необходимость установки дополнительных системных библиотек.

В контексте игрового движка ImageSharp используется для загрузки текстурных файлов, их декодирования и преобразования в формат, пригодный для передачи в OpenGL. Библиотека предоставляет простой и интуитивно понятный API, позволяющий загрузить изображение вызовом метода Image.Load, после чего получить доступ к пиксельным данным. Полученный массив пиксельных данных может быть передан в OpenGL через вызов функции glTexImage2D для создания текстурного объекта. Благодаря использованию механизмов SIMD, ImageSharp демонстрирует высокую производительность, приближающуюся к нативным решениям, что особенно важно при загрузке большого количества текстур при старте сцены.

При загрузке изображений, размер которых не превышает определённого порогового значения, библиотека гарантирует непрерывность пиксельных данных в памяти. Эта особенность позволяет передавать указатель на данные непосредственно в OpenGL без дополнительного копирования, что повышает эффективность работы с видеопамятью. Библиотека распространяется под лицензией Apache 2.0, что делает её свободно используемой в некоммерческих и образовательных проектах. Таким образом, выбор SixLabors.ImageSharp обусловлен её функциональностью, производительностью и отсутствием нативных зависимостей, что полностью соответствует требованиям разрабатываемого игрового движка.

\subsection{Общая структура движка}
Игровой движок состоит из следующих основных модулей: графической подсистемы на базе OpenGL, системы управления игровыми объектами и компонентами, менеджера сцен, системы камер и освещения, редактора сцен на Windows Forms, системы импорта моделей (STL, GLB), подсистемы обработки ввода и вспомогательных утилит. Графическая подсистема взаимодействует со сценой для получения отображаемых объектов. Редактор сцен через менеджер игровых объектов сцены управляет составом сцены и через графическую подсистему отображает трёхмерный предпросмотр. Система импорта моделей создает объекты и добавляет их в сцену через менеджер объектов. На рисунке \ref{fig:architecture} представлена общая схема архитектуры программной системы.

\begin{figure}[H]
	\centering
	\includegraphics[width=1.0\textwidth]{"images/Схема архитектуры.png"}
	\caption{Общая схема архитектуры приложения}
	\label{fig:architecture}
\end{figure}

Данная программная система содержит множество компонентов, взаимодействующих между собой. В таблице \ref{table21:table} представлены основные компоненты системы и их назначение.

\begin{xltabular}{\textwidth}{|>{\raggedright\arraybackslash}p{6cm}|>{\raggedright\arraybackslash}X|}
	\caption{Назначение компонентов программной системы\label{table21:table}}\\ \hline
	~ \centrow Имя & \centrow Описание \\ \hline
	\endfirsthead
	\continuecaption{Продолжение таблицы \ref{table21:table}}
	~ \centrow Имя & \centrow Описание \\ \hline 
	\finishhead
	Camera & Данный класс представляет собой камеру. \\ \hline
	FreeCameraController & Контроллер камеры. Поддерживает движение камеры и вращение камеры в углах yaw (рысканье) и pitch (тангаж). \\ \hline
	Material & Материал, который создается при загрузки моделей. \\ \hline
	Mesh & Меш модели. Выступает контейнером для примитивов и содержит матрицу в случае .glb модели.\\ \hline
	Primitive & Примитив геометрии. Содержит буферы и шейдер. \\ \hline
	Texture & Текстура. Является составной частью материала. \\ \hline
	Transform & Компонент игрового объекта. Управляет его положением, поворотом и масштабом. \\ \hline
	DrawableObject & Базовый класс для игровых объектов. \\ \hline
	GLBGameObject & Игровой объект, созданный на основе .glb файла. \\ \hline
	STLGameObject & Игровой объект, созданный на основе .stl файла. \\ \hline
	DirectionalLight & Является направленным источником освещения. \\ \hline
	PointLight & Является точечным источником освещения. \\ \hline
	Spotlight & Является прожекторным источником освещения. \\ \hline
	CameraManager & Менеджер камер. Управляет камерами в сцене. \\ \hline
	GameObjectManager & Менеджер игровых объектов. Управляет игровыми объектами в сцене. \\ \hline
	LightManager & Менеджер света. Управляет источниками освещения в сцене. \\ \hline
	SceneManager & Менеджер сцен. Управляет загрузкой и выгрузкой сцен. \\ \hline
	GLSLMacrosBuilder & Вспомогательный класс для шейдера. Валидирует и добавляет в шейдер макросы на основе параметров рендеринга. \\ \hline
	MacrosCache & Хранит названия макросов, соотнесенные с настройками рендеринга. \\ \hline
	TextureActivator & Активирует текстуры для шейдера. Хранит кэш последних привязанных текстур. \\ \hline
	OffScreenRender & Необходим для работы визуального редактора. Рисует объекты в текстуры. \\ \hline
	RenderSystem & Система рендеринга. Отрисовывает объекты сцены. \\ \hline
	Shader & Шейдер. Хранится в примитиве. \\ \hline
	ShadersPool & Пул шейдеров. Необходим для кэширования шейдеров одинаковых материалов. \\ \hline
	Mathematics & Хранит в себе различные математические операции, используемые движком. \\ \hline
	Quaternion & Кватернион. Используется для расчета поворота объекта. \\ \hline
	BaseScene & Базовый класс сцены. \\ \hline
	GameApplication & Главный класс движка. Управляет запуском приложения. \\ \hline
	SystemCalls & Используется движком для вывода диагностических сообщений в консоль. \\ \hline
	GLBImporter & Импортер .glb моделей. \\ \hline
	GLBModel & GLB модель. Используется для создания GLBGameObject. \\ \hline
	GLBScene & Контейнер для GLB моделей. \\ \hline
	Node & Узел иерархии узлов GLB файла. \\ \hline
	STLImporter & Импортер .stl моделей. \\ \hline
	STLModel & STL модель. Используется для создания STLGameObject.
\end{xltabular}

\subsection{Взаимодействие основных модулей}
Основные модули движка взаимодействуют следующим образом. Редактор сцен инициирует операции (импорт модели, сохранение сцены, настройка параметров игрового объекта), обращаясь к менеджеру сцен и графической подсистеме. Менеджер сцен управляет составом сцены и жизненным циклом объектов. Графическая подсистема запрашивает данные о сцене у менеджера сцен и выполняет рендеринг. Подсистема ввода передаёт события через сцену к игровым объектам. Диаграмма последовательности на рисунке \ref{fig:sequence_diagram} демонстрирует порядок вызовов при рендеринге игрового объекта.

\begin{figure}[H]
	\centering
	\includegraphics[width=1.0\textwidth]{"images/Диаграмма последовательности для рендера.png"}
	\caption{Диаграмма последовательности для цикла рендеринга объекта}
	\label{fig:sequence_diagram}
\end{figure}

\subsection{Графическая подсистема}
\subsubsection{Управление буферами вершин и индексов}
Для хранения геометрических данных в движке используются вершинные буферы (VBO) и индексные буферы (EBO). Вершинный буфер хранит массив вершин, каждая из которых содержит позицию (три координаты), нормаль (три компонента) и текстурные координаты (две координаты). Индексный буфер хранит последовательность индексов, определяющих порядок соединения вершин в треугольники. Все буферы инкапсулированы в классе Primitive. При создании примитива в игровом объекте выполняются следующие операции: генерация идентификаторов буферов; привязка буфера к целевой точке; загрузка данных в буфер; настройка атрибутов вершин и активация атрибутов через. Для организации атрибутов используется объект вершинного массива (VAO), который сохраняет все настройки атрибутов и привязки буферов. Это позволяет переключаться между разными моделями. При уничтожении меша буферы удаляются через glDeleteBuffers, что гарантирует освобождение памяти видеокарты.

\subsubsection{Компиляция и загрузка шейдерных программ}
В разработанном движке шейдерные программы собираются динамически на основе набора макросов, определяющих характеристики материала, модель освещения и глобальные настройки рендеринга. Класс Shader инкапсулирует процесс создания шейдерной программы: создаёт идентификатор программы, компилирует вершинный и фрагментный шейдеры, линкует их и сохраняет соответствие имён uniform-переменных их местоположениям.

Перед компиляцией шейдера класс GLSLMacrosBuilder формирует строку с макросами, которая вставляется в начало исходного кода шейдера. Для этого используются атрибуты GLSLMacros, которыми помечены элементы перечислений RenderSettings и ShadingModels. Это позволяет включать в шейдер только те участки кода, которые необходимы для текущего материала: определённые текстуры, факторы материалов или модели освещения. Глобальные настройки рендеринга (максимальная интенсивность бликов, сила окружающего освещения, использование монохромного окружающего света) передаются в шейдер через макросы, сформированные на основе класса GlobalRenderSettings. При необходимости в шейдер также включается макрос наличия матрицы меша.

После компиляции шейдер загружается в видеопамять. В случае ошибки компиляции или линковки движок выводит в консоль лог ошибок. При включённом режиме DEBUG\_MODE скомпилированные шейдеры сохраняются в текстовые файлы для отладки. В таблицах \ref{tab:render_settings} и \ref{tab:shading_models} представлены поддерживаемые флаги рендеринга и модели освещения, а в таблице \ref{tab:global_settings} — глобальные настройки.

\begin{xltabular}{\textwidth}{|>{\raggedright\arraybackslash}p{6cm}|>{\raggedright\arraybackslash}X|}
	\caption{Render settings\label{tab:render_settings}}\\ \hline
	~ \centrow Имя & \centrow Описание \\ \hline
	\endfirsthead
	\continuecaption{Продолжение таблицы \ref{tab:render_settings}}
	~ \centrow Имя & \centrow Описание \\ \hline 
	\finishhead
	VertexColors & Цвета вершин. \\
	\hline
	AlbedoMap & Основная текстура материала. \\
	\hline
	MetallicRoughnessMap & Карта металличности/шероховатости. \\
	\hline
	NormalMap & Карта нормалей. \\
	\hline
	OcclusionMap & Карта оклюзий. \\
	\hline
	EmissiveMap & Карта эмиссии. \\
	\hline
	BaseColorFactor & Общий цвет полигонов. \\
	\hline
	MetallicFactor & Значение металличности. \\
	\hline
	RoughnessFactor & Значение шероховатости. \\
	\hline
	EmissiveFactor & Значение эмиссии.
\end{xltabular}

\begin{xltabular}{\textwidth}{|>{\raggedright\arraybackslash}p{6cm}|>{\raggedright\arraybackslash}X|}
	\caption{Модели освещения\label{tab:shading_models}}\\ \hline
	~ \centrow Имя & \centrow Описание \\ \hline
	\endfirsthead
	\continuecaption{Продолжение таблицы \ref{tab:shading_models}}
	~ \centrow Имя & \centrow Описание \\ \hline 
	\finishhead
	BlinnPhong & Модель освещения Блинна-Фонга. \\
	\hline
	Gouraud & Модель освещения Гуро. \\
	\hline
	PBR MetallicRoughness & Модель освещения для physically based rendering материалов. Не используется. \\
	\hline
	Unlit & Отсутствие освещения. \\
\end{xltabular}

\begin{xltabular}{\textwidth}{|>{\raggedright\arraybackslash}p{6cm}|>{\raggedright\arraybackslash}X|}
	\caption{Модели освещения\label{tab:global_settings}}\\ \hline
	~ \centrow Имя & \centrow Описание \\ \hline
	\endfirsthead
	\continuecaption{Продолжение таблицы \ref{tab:global_settings}}
	~ \centrow Имя & \centrow Описание \\ \hline 
	\finishhead
	MaxShininess & Максимальная интенсивность зеркальных бликов (значение по умолчанию — 32) \\
	\hline
	AmbientStrength & Сила окружающего освещения (значение по умолчанию — 0.1) \\
	\hline
	UseMonochromeAmbient & Флаг использования монохромного окружающего освещения вместо цветного
\end{xltabular}

\subsection{Система трансформации}
Система трансформации отвечает за управление положением, поворотом и масштабом игровых объектов в трёхмерном пространстве. Каждый игровой объект содержит компонент трансформации, который хранит его параметры в мировом пространстве.

\subsubsection{Матрицы преобразований}
Для преобразования вершин объекта из локальных координат в мировые используется матрица модели (model matrix). Эта матрица получается путём последовательного перемножения матриц масштабирования, поворота и перемещения. Формула вычисления матрицы модели имеет следующий вид:

$$
\mathbf{M} = \mathbf{T} \times \mathbf{R} \times \mathbf{S}
$$

где $\mathbf{M}$ — итоговая матрица модели, $\mathbf{T}$ — матрица перемещения, $\mathbf{R}$ — матрица поворота, $\mathbf{S}$ — матрица масштабирования.

Матрица перемещения имеет следующий вид:

$$
\mathbf{T} = \begin{pmatrix}
	1 & 0 & 0 & t_x \\
	0 & 1 & 0 & t_y \\
	0 & 0 & 1 & t_z \\
	0 & 0 & 0 & 1
\end{pmatrix}
$$

где $t_x$, $t_y$, $t_z$ — координаты перемещения объекта в мировом пространстве.

Матрица масштабирования имеет следующий вид:

$$
\mathbf{S} = \begin{pmatrix}
	s_x & 0 & 0 & 0 \\
	0 & s_y & 0 & 0 \\
	0 & 0 & s_z & 0 \\
	0 & 0 & 0 & 1
\end{pmatrix}
$$

где $s_x$, $s_y$, $s_z$ — коэффициенты масштабирования по осям X, Y и Z соответственно.

В движке матрицы хранятся в виде указателей на массив чисел с плавающей точкой (float*), что позволяет эффективно передавать их в OpenGL без дополнительного копирования данных. Все матричные операции — умножение, транспонирование, вычисление обратной матрицы — выполняются статическими методами класса Mathematics. Для отображения объекта на экране матрица модели последовательно умножается на матрицу вида камеры и матрицу проекции, формируя итоговую матрицу MVP:

$$
\mathbf{MVP} = \mathbf{P} \times \mathbf{V} \times \mathbf{M}
$$

где $\mathbf{P}$ — матрица проекции, $\mathbf{V}$ — матрица вида, $\mathbf{M}$ — матрица модели. Эта матрица передаётся в вершинный шейдер и применяется к каждой вершине объекта.

\subsubsection{Позиция, поворот и масштаб}
Компонент трансформации хранит три основных параметра в мировом пространстве: позицию, поворот и масштаб. Позиция задаётся вектором из трёх координат, определяющим смещение объекта в пространстве. Поворот задаётся кватернионом — математическим объектом, представляющим вращение в трёхмерном пространстве. Масштаб задаётся вектором из трёх коэффициентов, позволяя независимо масштабировать объект по каждой оси.

Кватернион в движке реализован собственным классом и представляет собой четырёхкомпонентный вектор вида:

$$
\mathbf{q} = (x, y, z, w) = \sin\left(\frac{\theta}{2}\right) \cdot \mathbf{u} + \cos\left(\frac{\theta}{2}\right)
$$

где $\theta$ — угол поворота, $\mathbf{u}$ — единичный вектор оси вращения. Использование кватернионов вместо углов Эйлера позволяет избежать явления блокировки шарнирного соединения и обеспечивает плавную интерполяцию вращений.

Преобразование кватерниона в матрицу поворота выполняется по следующей формуле:

$$
\mathbf{R} = \begin{pmatrix}
	1 - 2y^2 - 2z^2 & 2xy - 2zw & 2xz + 2yw & 0 \\
	2xy + 2zw & 1 - 2x^2 - 2z^2 & 2yz - 2xw & 0 \\
	2xz - 2yw & 2yz + 2xw & 1 - 2x^2 - 2y^2 & 0 \\
	0 & 0 & 0 & 1
\end{pmatrix}
$$

где $(x, y, z, w)$ — компоненты кватерниона.

Компонент трансформации предоставляет поля для хранения позиции (Vector3), поворота (Quaternion) и масштаба (Vector3). При изменении любого из этих параметров пересчитывается матрица модели через вызов методов класса Mathematics. Матрица модели хранится в виде массива из шестнадцати чисел с плавающей точкой, готового к передаче в OpenGL. Для удобства работы с поворотом также доступны методы для установки вращения через углы Эйлера, которые внутренне преобразуются в кватернион с использованием соответствующих математических функций.

\subsection{Система камер}
Система камер отвечает за формирование изображения сцены с определённой точки обзора. В разработанном движке камера реализована классом Camera, который хранит параметры проекции, положение и ориентацию в пространстве.

\subsubsection{Типы проекций (перспективная и ортографическая)}
В движке поддерживаются два типа проекций. Перспективная проекция используется в игровом режиме и характеризуется углом обзора (FOV) и соотношением сторон экрана.

\subsubsection{Управление положением и направлением обзора}
Положение камеры задаётся вектором Position. Направление обзора определяется тремя векторами: Front (направление взгляда), RightVector (правая сторона) и Up (верх). Управление направлением осуществляется через углы рыскания (Yaw) и тангажа (Pitch). При их изменении вызывается метод CalculateVectors, который пересчитывает векторы Front, RightVector и Up.

\subsubsection{Матрица вида и матрица проекции}
Матрица вида (View Matrix) преобразует мировые координаты в координаты пространства камеры. В движке она вычисляется как произведение матрицы ориентации и матрицы перемещения.

Матрица ориентации строится из векторов RightVector, Up и Front:

$$
\mathbf{R} = \begin{pmatrix}
	r_x & r_y & r_z & 0 \\
	u_x & u_y & u_z & 0 \\
	-f_x & -f_y & -f_z & 0 \\
	0 & 0 & 0 & 1
\end{pmatrix}
$$

где $(r_x, r_y, r_z)$ — компоненты RightVector, $(u_x, u_y, u_z)$ — компоненты Up, $(f_x, f_y, f_z)$ — компоненты Front.

Матрица перемещения строится из позиции камеры:

$$
\mathbf{T} = \begin{pmatrix}
	1 & 0 & 0 & -p_x \\
	0 & 1 & 0 & -p_y \\
	0 & 0 & 1 & -p_z \\
	0 & 0 & 0 & 1
\end{pmatrix}
$$

где $(p_x, p_y, p_z)$ — компоненты Position.

Итоговая матрица вида вычисляется как произведение:

$$
\mathbf{V} = \mathbf{R} \times \mathbf{T}
$$

Матрица проекции (Projection Matrix) преобразует координаты из пространства камеры в экранные координаты. Для перспективной проекции матрица имеет следующий вид:

$$
\mathbf{P}_{persp} = \begin{pmatrix}
	\frac{1}{aspect \cdot \tan(fov/2)} & 0 & 0 & 0 \\
	0 & \frac{1}{\tan(fov/2)} & 0 & 0 \\
	0 & 0 & -\frac{Far + Near}{Far - Near} & -\frac{2 \cdot Far \cdot Near}{Far - Near} \\
	0 & 0 & -1 & 0
\end{pmatrix}
$$

Для ортографической проекции матрица имеет следующий вид:

$$
\mathbf{P}_{ortho} = \begin{pmatrix}
	\frac{2}{right - left} & 0 & 0 & -\frac{right + left}{right - left} \\
	0 & \frac{2}{top - bottom} & 0 & -\frac{top + bottom}{top - bottom} \\
	0 & 0 & -\frac{2}{Far - Near} & -\frac{Far + Near}{Far - Near} \\
	0 & 0 & 0 & 1
\end{pmatrix}
$$

В движке обе матрицы хранятся в виде указателей на массивы из шестнадцати чисел с плавающей точкой в неуправляемой памяти. Матрица вида заполняется через вызов Mathematics.MultiplyMatrices4. Матрица проекции заполняется напрямую в зависимости от типа камеры.

\subsection{Система освещения}
\subsubsection{Типы источников света (точечные, направленные, окружающие)}
В разработанном игровом движке реализована система освещения, поддерживающая три типа источников света, каждый из которых имеет свои особенности и область применения.

Окружающее освещение (Ambient Light) представляет собой базовое равномерное освещение всей сцены, не зависящее от положения источников и направления лучей. Этот тип освещения имитирует рассеянный свет, многократно отражённый от поверхностей объектов и стен. Окружающее освещение не создаёт теней и не подчёркивает объём объектов, но обеспечивает видимость всех элементов сцены даже в отсутствие прямых источников. В движке окружающее освещение задаётся двумя параметрами: параметром — AmbientStrength, который задает интенсивность освещения и параметром UseMonochromeAmbient, который задает цвет амбиентного освещения. Если он выставлен в true, то освещение будет серым вне зависимости от цвета объекта. в ином случае данное освещение будет цвета объекта. Окружающее освещение применяется ко всем объектам сцены глобально.

Точечный источник света (Point Light) излучает свет равномерно во все стороны из одной точки в пространстве, подобно обычной лампочке. Такой источник характеризуется положением в мировых координатах, цветом, интенсивностью и радиусом действия, который задается параметрами Constant, Linear и Quadratic. Интенсивность освещения от точечного источника убывает с расстоянием по квадратичному закону, что создаёт естественный эффект затухания. В движке точечные источники могут быть добавлены в сцену в любом количестве, каждый из них независимо освещает окружающие объекты. Точечные источники подходят для имитации ламп, факелов, взрывов и других локализованных источников света.

Направленный источник света (Directional Light) имитирует бесконечно удалённый источник, такой как Солнце. В отличие от точечного, направленный источник не имеет позиции в пространстве, а характеризуется только направлением лучей. Все лучи направленного источника параллельны друг другу, поэтому освещённость объекта не зависит от расстояния до источника. В движке направленный источник задаётся углами Эйлера Yaw и Pitch, цветом и интенсивностью. Этот тип освещения используется для создания дневного освещения на открытых пространствах.

Для всех типов источников света в шейдере реализована модель освещения Фонга, включающая вычисление рассеянной (diffuse) и зеркальной (specular) составляющих для точечных и направленных источников. Окружающее освещение добавляется на финальном этапе как константная составляющая, не зависящая от геометрии сцены.

Прожектор (Spotlight) Четвёртым типом источника света, реализованным в движке, является прожектор. Прожектор представляет собой точечный источник, излучающий свет не во все стороны, а в пределах заданного конуса. Такой тип освещения имитирует направленные источники света: фонарик, автомобильную фару, уличный прожектор или луч света на сцене.

Прожектор характеризуется следующими параметрами: положение в мировых координатах, определяющее точку испускания света; Углы Эйлера Yaw и Pitch, задающее ось конуса; цвет и интенсивность излучения; угол внутреннего конуса, внутри которого свет достигает максимальной яркости; угол внешнего конуса, за пределами которого освещение полностью отсутствует; а также радиус действия, ограничивающий максимальное расстояние распространения света.

\subsubsection{Модель освещения}
В разработанном игровом движке реализованы две модели освещения: модель Гуро и модель Блинна-Фонга, переключение между которыми выполняется через параметр шейдера.

Модель освещения Гуро вычисляет интенсивность освещения в вершинах многоугольника с последующей линейной интерполяцией полученных значений по всей поверхности грани. В вершинном шейдере для каждой вершины рассчитываются цветовые компоненты: окружающая, рассеянная и зеркальная. Полученные цвета интерполируются растеризатором, и для каждого фрагмента используется уже вычисленное значение. Основным преимуществом модели Гуро является высокая производительность, поскольку вычисления освещения выполняются только для вершин, а не для каждого пикселя. Недостатком является заметная ступенчатость освещения на крупных полигонах, особенно в области зеркальных бликов.

Модель освещения Блинна-Фонга вычисляет освещение для каждого фрагмента независимо. Во фрагментный шейдер передаются нормаль, позиция и текстурные координаты, интерполированные из вершинных атрибутов. Для каждого фрагмента рассчитываются три составляющие итогового цвета.

Формула модели Блинна-Фонга имеет следующий вид:

$$
I = I_{ambient} + I_{diffuse} + I_{specular}
$$

где $I_{ambient}$ — компонента окружающего освещения:

$$
I_{ambient} = k_a \cdot I_a
$$

Компонента рассеянного освещения (диффузная):

$$
I_{diffuse} = k_d \cdot I_d \cdot \max(0, \vec{L} \cdot \vec{N})
$$

где $k_d$ — коэффициент диффузного отражения материала, $I_d$ — интенсивность источника света, $\vec{L}$ — единичный вектор от точки поверхности к источнику света, $\vec{N}$ — нормаль к поверхности.

Компонента зеркального отражения по модели Блинна-Фонга:

$$
I_{specular} = k_s \cdot I_s \cdot (\vec{N} \cdot \vec{H})^{n}
$$

где $k_s$ — коэффициент зеркального отражения, $I_s$ — интенсивность зеркальной составляющей источника, $\vec{H}$ — половинный вектор, вычисляемый как $\vec{H} = \frac{\vec{L} + \vec{V}}{|\vec{L} + \vec{V}|}$, $\vec{V}$ — вектор направления на камеру, $n$ — показатель степени, определяющий размер блика (коэффициент шероховатости поверхности).

Основным отличием модели Блинна-Фонга от классической модели Фонга является использование половинного вектора $\vec{H}$ вместо вычисления вектора отражения, что упрощает вычисления и даёт более реалистичные блики.

В разработанном движке модель освещения поддерживает до 8 источников света одновременно, каждый из которых может быть точечным, направленным или прожектором. Окружающее освещение задаётся глобально для всей сцены. Все коэффициенты освещения (диффузный, зеркальный, интенсивность) передаются в шейдер через uniform-переменные и могут изменяться в реальном времени через редактор материалов.

\subsubsection{Работа с SSBO}
Для эффективной передачи данных о источниках света в шейдерную программу в движке используется механизм Shader Storage Buffer Object. SSBO представляет собой буфер, доступный для чтения и записи как со стороны графического процессора, так и со стороны центрального процессора, что позволяет хранить в видеопамяти структурированные массивы данных произвольного размера. На рисунке \ref{fig:ssbo_structure} показана схема хранения данных света в SSBO на примере направленного источника света.

\begin{figure}[H]
	\centering
	\includegraphics[width=1.0\textwidth]{"images/Схема хранения света в ssbo.png"}
	\caption{Схема хранения данных в памяти}
	\label{fig:ssbo_structure}
\end{figure}

В движке реализован обобщённый класс ShaderStorageBufferManager<T>, параметризуемый типом структуры, описывающей данные источника света. Данный класс инкапсулирует все операции по управлению буфером: выделение памяти, добавление, обновление и удаление элементов, а также автоматическое расширение буфера при превышении начальной ёмкости.

При инициализации менеджера в конструктор передаются начальная ёмкость буфера и точка привязки (binding point). Внутри конструктора выполняется выделение памяти в оперативной памяти через NativeMemory.Alloc и в видеопамяти через вызов glGenBuffer и glBufferData. Для хранения количества активных источников в начале буфера отводится дополнительное пространство размером PADDING, что позволяет шейдеру знать, сколько элементов обрабатывать. Указатель на выделенную память хранится в поле \_lightData, а идентификатор буфера OpenGL — в свойстве LightHandler.

Метод AddLight добавляет новый источник света в буфер. При добавлении проверяется, не превышает ли текущее количество активных источников заявленную ёмкость. При превышении ёмкости буфер автоматически расширяется в два раза через вызов NativeMemory.Realloc, после чего обновляются данные в видеопамяти. При наличии свободных мест (отслеживаемых через поле \_freeSpaces) новый источник размещается на позиции ранее удалённого элемента. Это позволяет избежать фрагментации данных и пересортировки массива. После добавления обновляется счётчик активных источников в начале буфера.

Метод RemoveLight выполняет логическое удаление источника света, устанавливая поле isActive в соответствующей структуре в ноль. Счётчик активных источников уменьшается, а счётчик свободных мест увеличивается. Физического удаления данных из буфера не происходит, что позволяет избежать дорогостоящего сдвига массива и сохраняет производительность при частых удалениях.

Метод UpdateLight обновляет данные существующего источника по указанному индексу. Вычисляется смещение в буфере как произведение размера структуры на индекс плюс отступ PADDING, после чего через Marshal.StructureToPtr данные копируются в управляемую память, а затем через glBufferSubData — в видеопамять.

Метод Clear полностью очищает буфер, устанавливая счётчик активных источников в ноль и перезаписывая данные в видеопамяти через glBufferData.

Все операции с видеопамятью выполняются через привязку буфера к цели ShaderStorageBuffer с последующей отвязкой. Буфер также привязывается к указанной точке привязки через glBindBufferBase, что позволяет шейдеру обращаться к данным по этой точке.

Использование SSBO для хранения источников света даёт следующие преимущества: отсутствует ограничение на максимальное количество источников, характерное для uniform-массивов; данные могут динамически изменяться в рантайме без перекомпиляции шейдера; массив структур хранится в памяти GPU как непрерывный блок, что обеспечивает эффективный доступ из шейдера. В разработанном движке SSBO управляет массивами точечных источников, направленных источников и прожекторов, каждый из которых описывается отдельной структурой.

\subsection{Управление сценами}
\subsubsection{Загрузка и выгрузка сцен}
Менеджер сцен обеспечивает управление жизненным циклом игровых сцен, включая их загрузку, выгрузку, переключение и организацию истории переходов. Все загруженные сцены регистрируются в реестре по имени, что позволяет обращаться к ним без необходимости хранения прямых ссылок.

\textbf{Загрузка сцены}. При загрузке сцены менеджер сначала проверяет, не является ли запрашиваемая сцена уже текущей, чтобы избежать повторной инициализации. Если текущая сцена существует, для неё вызывается метод выхода, который приостанавливает игровую логику и уведомляет подписанные компоненты о завершении работы сцены. Затем выполняется очистка ресурсов текущей сцены через метод Cleanup, который освобождает все занятые объектом ресурсы: шейдерные программы, VBO, EBO, VAO, текстурные объекты, источники света и камеры. После очистки новая сцена устанавливается как текущая, ей передаётся ссылка на менеджер сцен, после чего выполняется десериализация сцены из файла, восстанавливающая все игровые объекты, их компоненты и параметры. Если приложение находится не в режиме редактирования, для новой сцены вызывается метод старта, инициирующий игровую логику. По завершении загрузки генерируются события, оповещающие другие подсистемы движка об изменении текущей сцены.

\textbf{Выгрузка сцены}. При выгрузке текущая сцена удаляется из памяти. Если приложение не в режиме редактирования, сначала вызывается метод выхода, дающий сцене возможность корректно завершить игровые процессы. Затем выполняется метод Cleanup, который рекурсивно обходит все игровые объекты сцены и для каждого из них освобождает связанные ресурсы: удаляются шейдерные программы, уничтожаются вершинные и индексные буферы через вызов соответствующих функций OpenGL, освобождаются текстурные объекты, очищаются массивы источников света в SSBO-буферах, а также освобождаются указатели камер и других компонентов. После завершения очистки ссылка на текущую сцену обнуляется, и генерируются события об изменении состояния. Менеджер сцен не удаляет саму сцену из реестра, что позволяет загрузить её повторно в будущем без повторной регистрации.

\textbf{Стек сцен}. Менеджер поддерживает стек сцен, позволяющий временно переключаться между сценами с возможностью возврата. Метод PushScene сохраняет текущую сцену в стеке и загружает новую. Метод PopScene выгружает текущую сцену и восстанавливает сцену из вершины стека. Этот механизм полезен для реализации наложения экранов, например, при открытии меню паузы поверх игровой сцены.

\textbf{Пауза сцены}. Менеджер также управляет состоянием паузы текущей сцены. При установке флага паузы вызываются соответствующие методы сцены, приостанавливающие обновление игровой логики без выгрузки ресурсов. Возобновление работы сцены происходит при снятии паузы.

\textbf{Полная очистка}. При завершении работы приложения или при необходимости полного сброса состояния вызывается метод Cleanup, который очищает все сцены в стеке, реестре и текущую сцену, а также отписывает все события. Этот метод гарантирует, что все неуправляемые ресурсы OpenGL будут освобождены до закрытия приложения, предотвращая утечки памяти.

\subsubsection{Сохранение сцены в файл и восстановление}
Для обеспечения возможности сохранения и последующей загрузки игровых сцен в движке реализован механизм сериализации в JSON-формат. Каждая сцена сохраняется в отдельный файл, имя которого соответствует имени класса сцены, что позволяет однозначно идентифицировать сцену при загрузке.

\textbf{Формат файла сцены}. Файл сцены имеет JSON-структуру, включающую два основных массива: массив путей к файлам моделей и массив описаний игровых объектов. Пример структуры сохранённой сцены представлен ниже: \\ \\

\begin{lstlisting}[caption={Пример JSON-файла сохранённой сцены}, label=lst:scene_json]
	{
		"Models": [
		"C:\\Users\\it_ge\\Desktop\\Amethyst-game-engine\\Client\\bin\\Debug\\net10.0-windows10.0.17763.0\\Models\\Room for debug.glb"
		],
		"GameObjects": [
		{
			"Tag": "Untagged",
			"Type": "GLB",
			"ModelIndex": 0,
			"GLBModelIndex": 0,
			"Transform": {
				"Position": [0, 0, 0],
				"Rotation": [0, 0, 0],
				"Scale": [1, 1, 1]
			}
		}
		]
	}
\end{lstlisting}

Массив путей к моделям содержит абсолютные пути к файлам импортированных моделей в форматах STL и GLB. Это позволяет при десериализации однократно загрузить каждую модель и использовать её для создания нескольких игровых объектов без дублирования данных. Массив игровых объектов содержит для каждого объекта следующие параметры: строковый тег для идентификации объекта, тип модели (STL или GLB), индекс модели в массиве путей, для GLB-моделей также индекс модели внутри составного файла, а также объект трансформации, включающий позицию, поворот и масштаб по трём осям. Позиция, поворот и масштаб хранятся в виде массивов из трёх чисел с плавающей точкой.

\textbf{Сериализация сцены}. Процесс сохранения сцены инициируется вызовом метода SerializeScene. В начале работы метода проверяется наличие папки Saves в директории приложения; при отсутствии папка создаётся автоматически. Затем формируется JSON-объект сцены: создаются пустые массивы для моделей и игровых объектов. Менеджер игровых объектов обходится для сбора всех используемых моделей, пути к которым добавляются в массив моделей. Далее каждый игровой объект преобразуется в JSON-объект с заполнением его тега, типа, индекса модели и параметров трансформации. Для объектов, импортированных из GLB, дополнительно сохраняется индекс подсцены внутри файла. Сформированный JSON-объект записывается в файл с отступами для удобства чтения и отладки.

\textbf{Десериализация сцены}. Восстановление сцены выполняется методом DeserializeScene. Сначала определяется полный путь к файлу сцены на основе имени текущего класса сцены. Если файл не существует, загрузка не выполняется и сцена остаётся пустой. При наличии файла происходит чтение и парсинг JSON-документа. Из корневого элемента извлекается массив путей к моделям, и для каждого пути определяется формат файла (по расширению). Каждая модель загружается через соответствующий импортёр (GLBImporter или STLImporter) и сохраняется во временном списке. Затем для каждого элемента из массива игровых объектов извлекаются индекс модели, тег и параметры трансформации. В зависимости от типа модели создаётся экземпляр соответствующего класса игрового объекта (GLBGameObject или STLGameObject). Созданному объекту присваивается тег, а затем через менеджер игровых объектов он добавляется на сцену. После добавления объекту устанавливаются параметры трансформации: позиция, поворот и масштаб, извлечённые из JSON.

\textbf{Управление ресурсами при загрузке}. Важной особенностью реализации является то, что модели загружаются один раз для всей сцены и затем используются многократно. Это позволяет экономить видеопамять и избегать дублирования геометрических данных. При удалении сцены или завершении приложения все загруженные модели и связанные с ними ресурсы OpenGL освобождаются через механизм очистки.

\textbf{Обработка ошибок}. При десериализации сцены предусмотрена обработка ошибочных ситуаций: если файл модели по указанному пути не существует, в консоль выводится сообщение об ошибке и программа завершается. При отсутствии файла сцены в целом процесс загрузки не выполняется, и сцена инициализируется в пустом состоянии, что позволяет избежать аварийного завершения приложения.

\subsection{Подсистема обработки ввода}
\subsubsection{Работа хуков}
Под хуками в контексте движка понимаются методы-обработчики, которые вызываются в ответ на определённые события жизненного цикла или действия пользователя. Все хуки определены как виртуальные методы в базовом классе сцены, что позволяет разработчику переопределять их в конкретных сценах для реализации специфической логики.

\textbf{Источники вызовов}. Хуки вызываются из двух основных мест: из класса GameApplication и из менеджера сцен. GameApplication является точкой входа приложения и получает события от операционной системы через колбэки GLFW. При возникновении события (нажатие клавиши, движение мыши, обновление кадра) GameApplication вызывает соответствующий метод у текущей сцены. Менеджер сцен, в свою очередь, вызывает хуки жизненного цикла (старт, выход, пауза, возобновление) в моменты переключения сцен или изменения состояния приложения.

\textbf{Жизненный цикл сцены}. При первой загрузке сцены менеджер сцен вызывает хук OnStart. Этот метод предназначен для однократной инициализации логики сцены, создания начальных объектов и подписки на события. При выгрузке сцены вызывается хук OnExit, который даёт возможность выполнить завершающие действия: сохранить состояние, освободить временные ресурсы и т.д. При приостановке сцены (например, при открытии меню паузы) вызывается хук OnPause, а при возобновлении — OnResume. Это позволяет корректно управлять игровым процессом без потери данных.

\textbf{Игровой цикл}. Основной игровой цикл реализован в классе GameApplication. На каждом кадре последовательно вызываются хуки Update и FixedUpdate. Хук Update вызывается с параметром deltaTime, представляющим время, прошедшее с предыдущего кадра. Этот хук предназначен для обработки логики, зависящей от частоты кадров: движения объектов, анимации, проверки условий. Хук FixedUpdate вызывается с фиксированным временным интервалом, что обеспечивает детерминированное поведение физических расчётов независимо от производительности оборудования.

\textbf{Обработка ввода}. При получении события от операционной системы GameApplication вызывает соответствующий хук у текущей сцены. Реализованы следующие хуки ввода: OnKeyDown и OnKeyUp для событий клавиатуры; OnMouseDown, OnMouseUp, OnMouseMove и OnMouseWheel для событий мыши. Каждый хук получает аргументы события, содержащие подробную информацию: код нажатой клавиши, модификаторы, координаты курсора, величину прокрутки колеса.

\textbf{Распространение событий на игровые объекты}. Ключевой особенностью реализации является то, что все перечисленные хуки в базовом классе сцены не содержат собственной логики, а лишь выполняют обход всех игровых объектов, зарегистрированных на сцене, и вызывают соответствующий метод у каждого объекта. Это позволяет любому игровому объекту переопределить нужные хуки и реагировать на события без необходимости регистрации в диспетчере событий или подписки на делегаты. Такой подход упрощает разработку игровой логики: разработчику достаточно создать класс, унаследованный от GameObject, и переопределить нужные виртуальные методы.

\textbf{Порядок вызовов}. Для каждого события гарантируется определённый порядок вызовов. В случае хуков жизненного цикла и игрового цикла все игровые объекты обрабатываются в порядке их добавления на сцену. Для событий ввода порядок не регламентируется, но гарантируется, что все объекты получат событие. Это позволяет реализовывать простые сценарии взаимодействия без необходимости управления очередью обработчиков.

\textbf{Особенности реализации в редакторе}. При работе в режиме редактирования хуки не вызываются, чтобы разработчик мог спокойно редактировать сцену. При переключении в режим игры работа игрового цикла возобновляется, и все хуки начинают вызываться в штатном режиме.

\subsubsection{Рейкастинг}
Рейкастинг представляет собой механизм определения объекта в трёхмерной сцене, на который указывает курсор мыши. В разработанном движке этот механизм используется для реализации интерактивного взаимодействия пользователя с объектами сцены, например, для выбора объекта кликом мыши.

\textbf{Принцип работы}. Рейкастинг основан на построении луча из камеры через позицию курсора на экране и проверке пересечения этого луча с ограничивающими объёмами объектов. Луч задаётся начальной точкой (позиция камеры) и направлением, которое вычисляется на основе координат курсора, размера окна и параметров проекции камеры. Направление луча определяется таким образом, чтобы он проходил через точку на ближней плоскости отсечения, соответствующую положению курсора, и уходил вглубь сцены.

\textbf{Построение луча}. Метод GetPickRay класса Camera принимает координаты курсора в оконной системе координат, ширину и высоту окна, а также матрицы проекции и вида. Вычисление выполняется в несколько этапов: преобразование оконных координат в нормализованные координаты устройства, затем в координаты пространства камеры и, наконец, в мировые координаты. Полученный луч используется для проверки пересечения с объектами сцены. Все математические операции выполняются с использованием матричных преобразований, что обеспечивает корректную работу при любых настройках камеры.

\textbf{Проверка пересечения}. Для каждого игрового объекта на сцене вызывается метод RayIntersectsAABB, который проверяет, пересекает ли луч ограничивающий прямоугольник объекта в мировых координатах. Использование ограничивающего прямоугольника, а не точной геометрии модели, является компромиссом между точностью и производительностью. Вычисление точного пересечения с каждым полигоном модели потребовало бы значительно больше вычислительных ресурсов и могло бы привести к снижению частоты кадров при большом количестве объектов.

\textbf{Выбор ближайшего объекта}. Поскольку луч может пересекать несколько объектов одновременно (например, когда один объект находится за другим), необходимо определить, какой из них расположен ближе к камере. Алгоритм перебирает все объекты сцены, для каждого вычисляет расстояние от начала луча до точки пересечения и сохраняет объект с минимальным положительным расстоянием. Расстояние считается положительным только в том случае, если точка пересечения находится перед камерой (а не позади неё). В результате возвращается объект, ближайший к камере среди всех пересечённых.

\textbf{Обработка результатов}. Метод PickObject возвращает объект типа DrawableObject, если пересечение найдено, или null в противном случае. Полученный объект может быть использован для дальнейшей обработки: подсветки, удаления, изменения свойств или выполнения игровой логики. В демонстрационном приложении результат рейкастинга используется для выбора объекта при клике мыши.

\textbf{Производительность}. В текущей реализации рейкастинг выполняется с линейной сложностью относительно количества объектов на сцене, так как перебираются все объекты. При сценах с небольшим количеством объектов (до нескольких сотен) этого достаточно для поддержания 60 кадров в секунду. Для сцен с тысячами объектов требуется дополнительная оптимизация, например, пространственное разбиение сцены, но в рамках дипломной работы такая оптимизация не реализована.

\subsection{Система импорта трёхмерных моделей}
Система импорта трёхмерных моделей отвечает за загрузку файлов внешних форматов, их парсинг и преобразование во внутреннее представление движка, пригодное для использования в рендеринге.

\subsubsection{Поддерживаемые форматы (STL, GLB)}
В разработанном игровом движке реализована поддержка двух форматов трёхмерных моделей: STL и GLB.

\textbf{Формат STL} (Stereolithography) является широко распространённым форматом для хранения трёхмерных моделей, особенно в области 3D-печати и инженерных расчётов. STL представляет геометрию объекта в виде набора треугольников, каждый из которых описывается тремя вершинами и вектором нормали. Формат существует в двух вариантах: текстовом (ASCII) и бинарном (Binary). Текстовый STL легко читается человеком и отлаживается, но занимает больше места на диске. Бинарный STL более компактен и быстрее загружается, поскольку не требует парсинга текстовых чисел. STL-файлы не поддерживают иерархию узлов, материалы, текстурные координаты и анимацию, что ограничивает их применение в сложных игровых проектах. Однако простота формата позволяет быстро импортировать модели, созданные в САПР-системах, что делает его удобным для прототипирования и демонстрационных целей.

\textbf{Формат GLB} является бинарной версией формата glTF (GL Transmission Format), разработанного консорциумом Khronos Group как современный стандарт для передачи трёхмерных сцен в веб-приложениях и игровых движках. В отличие от STL, GLB является полнофункциональным форматом, поддерживающим: иерархическую структуру узлов с матрицами трансформации; несколько сеток в одном файле; материалы с настраиваемыми параметрами (цвет, металличность, шероховатость); текстурные координаты и загрузку текстур; анимацию трансформаций; скелетную анимацию (включая кости и веса вершин). GLB представляет собой один бинарный файл, в котором объединены JSON-описание сцены и бинарные данные геометрии, что упрощает распространение моделей, так как все компоненты хранятся в одном файле, а не в наборе файлов. Этот формат выбран в качестве основного для сложных моделей, требующих поддержки материалов и иерархии. Устройство файла показано на рисунке 
\ref{fig:glb_structure}

\begin{figure}[H]
	\centering
	\includegraphics[width=1.0\textwidth]{"images/gltf файл.jpg"}
	\caption{Структура glb файла}
	\label{fig:glb_structure}
\end{figure}

\textbf{Выбор форматов}. Поддержка STL и GLB обусловлена их взаимодополняющими характеристиками. STL используется для быстрого импорта простых моделей и прототипирования, а также для совместимости с системами, экспортирующими только этот формат. GLB применяется для полноценных игровых моделей с материалами и анимацией. Такая комбинация покрывает широкий спектр задач от инженерных демонстраций до полноценных игр.

\subsubsection{Парсинг данных из STL-файлов}
Импорт STL-файлов в движке реализован с использованием бинарного чтения, что обеспечивает высокую производительность даже при загрузке моделей с большим количеством треугольников. Процесс парсинга начинается с проверки существования файла по указанному пути; при отсутствии файла в консоль выводится сообщение об ошибке, а загрузка прерывается.

Определение типа STL-файла выполняется путём чтения первых пяти байт. Если файл начинается с последовательности "solid", он идентифицируется как текстовый STL. В текущей версии движка текстовый формат не поддерживается, так как его обработка требует значительно больше вычислительных ресурсов и времени, а область применения таких моделей в игровой разработке ограничена. При обнаружении текстового формата в консоль выводится сообщение о неподдерживаемом формате, и загрузка прекращается. Все поддерживаемые файлы обрабатываются как бинарные STL.

После определения формата указатель чтения возвращается в начало файла, и считывается 80-байтовый заголовок, который обычно содержит произвольную текстовую информацию о модели и может быть проигнорирован. Затем извлекается 4-байтовое целое без знака, указывающее количество треугольников в модели. Перед загрузкой выполняется проверка: если количество треугольников превышает заданный лимит (один миллион), загрузка прерывается с выводом сообщения об ошибке, что предотвращает чрезмерное потребление памяти.

Для хранения геометрических данных создаётся структура, содержащая массивы вершин и нормалей. Вершины и нормали хранятся в виде плоских массивов байтов, что позволяет эффективно копировать данные из файла и впоследствии передавать их в OpenGL. Размер массивов вычисляется заранее на основе количества треугольников: каждый треугольник содержит три вершины, каждая вершина описывается тремя координатами с плавающей точкой, что даёт девять чисел с плавающей точкой на треугольник.

В цикле по всем треугольникам последовательно считываются данные: сначала нормаль треугольника (три числа с плавающей точкой), затем координаты трёх вершин (девять чисел с плавающей точкой). Нормаль записывается в массив нормалей трижды — по одному разу для каждой вершины треугольника, что упрощает последующее использование в шейдерах, где нормаль требуется для каждой вершины. Координаты вершин копируются в массив вершин.

После чтения вершин и нормалей считывается 2-байтовый атрибут, который в стандарте STL зарезервирован для пользовательских данных. В некоторых STL-файлах в этом поле кодируется цвет треугольника в формате RGB555 (5 бит на каждый цветовой канал). Движок поддерживает извлечение цветов из этого поля: выделяется 5 бит на красный, зелёный и синий каналы, каждый из которых преобразуется в значение с плавающей точкой в диапазоне от 0 до 1. Альфа-канал устанавливается в единицу (полностью непрозрачный). Если в файле обнаружены цветовые данные, для модели активируется режим вершинных цветов, который может быть использован в шейдере вместо или в дополнение к текстурам. Если цветовые данные отсутствуют, соответствующая опция рендеринга отключается.

После успешного чтения всех треугольников выполняется проверка, что указатель чтения достиг конца файла. Если после обработки всех треугольников в файле остались непрочитанные данные, файл считается повреждённым, и загрузка прерывается с сообщением об ошибке.

На завершающем этапе парсинга вычисляется ограничивающий прямоугольник модели. Для этого обходятся все загруженные вершины, и для каждой определяются минимальные и максимальные значения по каждой оси. Полученный прямоугольник будет использоваться для рейкастинга и отсечения невидимых объектов. На основе собранных данных создаётся объект модели, который сохраняет путь к исходному файлу для возможной перезагрузки и настройки рендеринга.

Все операции с файлом выполняются с использованием буферизованного чтения, что минимизирует количество обращений к диску. При возникновении любой ошибки в процессе чтения (неверный формат, преждевременный конец файла, повреждённые данные) загрузка прерывается, и в консоль выводится диагностическое сообщение, позволяющее определить причину сбоя.

\subsubsection{Парсинг данных и иерархии узлов из GLB-файлов}
Импорт GLB-файлов в разработанном движке реализован в соответствии со спецификацией glTF 2.0. GLB-файл представляет собой бинарный контейнер, объединяющий JSON-описание сцены и бинарные данные геометрии в одном файле. Процесс парсинга начинается с проверки существования файла и его размера: файлы размером более одного гигабайта не загружаются для предотвращения чрезмерного потребления памяти.

Первым этапом является чтение заголовка GLB-файла, который содержит магическое число (0x46546C67, ASCII-строка "glTF"), версию формата и размер файла. Движок поддерживает только версию 2.x, так как более старые версии формата не обеспечивают необходимой функциональности. При обнаружении неподдерживаемой версии загрузка прерывается с выводом сообщения об ошибке.

После заголовка последовательно считываются блоки данных (chunks). Первый блок является JSON-блоком, содержащим описание всей сцены в текстовом формате. Этот блок парсится в JSON-документ, который в дальнейшем используется для извлечения структуры сцены. Второй блок является бинарным блоком и содержит буферы с вершинными данными — позициями, нормалями, текстурными координатами и индексными данными.

Из JSON-дескриптора извлекаются массивы буферов, буферов представления (buffer views) и аксессоров (accessors). Буферы представляют собой сырые бинарные данные. Буферы представления определяют, какой участок буфера и с каким шагом между элементами должен использоваться. Аксессоры описывают, как интерпретировать данные из буфера представления: тип компонента (целое число со знаком/без знака, число с плавающей точкой), количество компонентов в элементе (скаляр, вектор из двух, трёх или четырёх компонентов), нормализованы ли значения, а также общее количество элементов. Аксессоры также могут содержать минимальные и максимальные значения координат, что позволяет вычислить ограничивающий прямоугольник модели без обхода всех вершин.

Особенностью реализации является поддержка разреженных аксессоров (sparse accessors), которые используются для эффективного хранения данных, где большинство значений являются дефолтными. При обнаружении разреженного аксессора движок копирует базовые данные, а затем заменяет элементы по указанным индексам значениями из разреженного хранилища. Это позволяет корректно загружать модели, оптимизированные с использованием данной техники.

Из JSON-дескриптора также извлекаются описания мешей, каждый из которых может содержать несколько примитивов. Примитив определяет тип геометрии (обычно треугольники), аксессор индексов и аксессоры атрибутов вершины: позиции (обязательный), нормали, текстурные координаты (до пяти наборов). Для каждого примитива может быть указан материал, определяющий внешний вид поверхности. На основе этих данных для каждого примитива создаётся структура, содержащая массивы вершин, индексов, нормалей, текстурных координат, а также ограничивающий прямоугольник.

Иерархия узлов (nodes) извлекается из JSON-дескриптора и реконструируется в граф сцены. Каждый узел может содержать дочерние узлы, матрицу трансформации или отдельные компоненты трансформации (позицию, поворот в виде кватерниона, масштаб). Для каждого узла вычисляется локальная матрица трансформации, а затем рекурсивно вычисляются глобальные матрицы с учётом трансформаций всех родительских узлов. Узлы, содержащие меши, преобразуются в игровые объекты с соответствующими компонентами трансформации и визуального отображения. Узлы, являющиеся костями скелетной анимации, идентифицируются и могут быть использованы для анимации (хотя сама анимация в текущей версии не поддерживается).

В таблице \ref{tab:glb_support} показано, то, какие возможности формата GLB поддерживает движок:

\begin{xltabular}{\textwidth}{|>{\raggedright\arraybackslash}p{6cm}|>{\raggedright\arraybackslash}X|}
	\caption{Поддерживаемые возможности формата GLB в разработанном импортёре\label{tab:glb_support}}\\ \hline
	~ \centrow Возможность & \centrow Описание поддержки \\ \hline
	\endfirsthead
	\continuecaption{Продолжение таблицы \ref{tab:glb_support}}
	~ \centrow Возможность & \centrow Описание поддержки \\ \hline 
	\finishhead
	
	Версия формата & Поддерживается только версия 2.x. При обнаружении других версий загрузка прерывается с выводом сообщения об ошибке. \\
	\hline
	Буферы (buffers) & Полная поддержка: чтение из бинарного блока GLB и из внешних URI в кодировке base64. \\
	\hline
	Буферы представления (buffer views) & Полная поддержка: извлечение данных с указанным смещением, длиной и шагом между элементами. \\
	\hline
	Аксессоры (accessors) & Полная поддержка: все типы компонентов (целые числа со знаком/без знака, числа с плавающей точкой), нормализация, все типы данных (скаляр, векторы). \\
	\hline
	Разреженные аксессоры (sparse accessors) & Поддерживаются. Выполняется замена элементов по указанным индексам. \\
	\hline
	Примитивы (primitives) & Поддерживаются все возможные типы. \\
	\hline
	Позиции вершин (POSITION) & Обязательный атрибут. Используется для построения геометрии и ограничивающего прямоугольника. \\
	\hline
	Нормали (NORMAL) & Поддерживаются. \\
	\hline
	Текстурные координаты (TEXCOORD\_0 — TEXCOORD\_4) & Поддерживаются до пяти наборов. \\
	\hline
	Материалы (materials) & Поддерживается модель PBR (Metallic-Roughness). \\
	\hline
	Базовый цвет (baseColorFactor) & Поддерживается. Задаётся в виде RGBA-вектора. \\
	\hline
	Карта базового цвета (baseColorTexture) & Поддерживается. \\
	\hline
	Металличность (metallicFactor) & Поддерживается. Значение от 0 до 1. \\
	\hline
	Шероховатость (roughnessFactor) & Поддерживается. Значение от 0 до 1. \\
	\hline
	Карта металличности и шероховатости (metallicRoughnessTexture) & Поддерживается. \\
	\hline
	Карта нормалей (normalTexture) & Поддерживается с учётом масштаба (scale). \\
	\hline
	Карта окклюзии (occlusionTexture) & Поддерживается с учётом силы эффекта (strength). \\
	\hline
	Карта эмиссии (emissiveTexture) & Поддерживается. \\
	\hline
	Фактор эмиссии (emissiveFactor) & Поддерживается. Задаётся в виде RGB-вектора. \\
	\hline
	Семплеры (samplers) & Поддерживаются параметры фильтрации (магнитная и минимальная) и режимы наложения по осям S и T (wrapS, wrapT). \\
	\hline
	Изображения (images) & Поддерживаются изображения, встроенные в буферы через base64. \\
	\hline
	Текстуры (textures) & Полная поддержка. Связывают изображение с семплером. \\
	\hline
	Узлы (nodes) & Полная поддержка иерархии, матриц трансформации и отдельных компонентов (позиция, поворот в виде кватерниона, масштаб). \\
	\hline
	Сцены (scenes) & Поддерживаются все сцены в файле. Каждая сцена импортируется как отдельная модель. \\
	\hline
	Скелетная анимация (skins, joints) & Идентифицируются узлы-кости, но сама анимация не воспроизводится. \\
	\hline
	Ограничение на размер файла & Максимальный размер файла — 1 гигабайт. При превышении загрузка прерывается. \\
	\hline
\end{xltabular}

После завершения парсинга всех компонентов для каждой сцены рекурсивно обходится граф узлов, и для каждого узла, содержащего меш, создаётся модель, которая объединяет все меши узла и всех его потомков. Каждая полученная модель сохраняет путь к исходному файлу и свой индекс в сцене, что позволяет в дальнейшем ссылаться на неё при сериализации сцены. Весь процесс импорта полностью управляем: при возникновении любой ошибки (неверный формат, отсутствие обязательных полей, повреждённые данные) загрузка прерывается, а в консоль выводится диагностическое сообщение.

\subsubsection{Выводы по разделу}
В рамках технического проекта определено внутреннее устройство разработанного игрового движка на языке C\# с использованием графического API OpenGL. Обоснован выбор C\#, OpenGL, Windows Forms, SixLabors.ImageSharp и библиотек для импорта STL и GLB в качестве технологической основы проекта. Описана архитектура движка, включающая графическую подсистему, систему управления игровыми объектами, менеджер сцен, систему камер и освещения, а также редактор сцен. Подробно рассмотрено устройство системы трансформации с использованием матриц преобразований и кватернионов, включая формулы вычисления матриц модели, вида и проекции. Описана графическая подсистема:управление вершинными и индексными буферами, компиляция шейдеров с динамической генерацией макросов на основе характеристик материала и модели освещения, а также конвейер рендеринга. Отдельно рассмотрена система импорта трёхмерных моделей из форматов STL и GLB с поддержкой иерархии узлов, материалов, текстур и разреженных аксессоров.

В качестве практического подтверждения проектных решений рассмотрено демонстрационное приложение — трёхмерная игра, в которой пользователь нажимает на объекты, после чего они исчезают и возрождаются в случайных позициях. Приложение демонстрирует совместную работу обработки ввода, рейкастинга для выбора объектов, системы трансформации и рендеринга, что подтверждает работоспособность разработанного движка.